\documentclass[../../main.tex]{subfiles}% 注意这里的文件路径不能用 ./main.tex ,否则用latexmk编译子文件会报错
\graphicspath{{\subfix{./image/}}} % 指定图片目录,后续可以直接使用图片文件名
% 注意这里的文件路径不能用 ../../image/ ,否则用latexmk编译子文件会报错

% 例如:
% \begin{figure}[H]
% \centering
% \includegraphics[scale=0.3]{图.png}
% \caption{}
% \label{figure:图}
% \end{figure}
% 注意:上述\label{}一定要放在\caption{}之后,否则引用图片序号会只会显示??.

\begin{document}

\section{常用初等不等式}

\begin{proposition}[指数/对数函数不等式]\label{proposition:指数对数函数不等式}
\begin{enumerate}[(1)]
\item\label{proposition:指数对数函数不等式1} $\ln \left( 1+x \right) <\frac{x}{\sqrt{1+x}},x>0\iff \ln x<\sqrt{x}-\frac{1}{\sqrt{x}},x>1.$

\item \label{proposition:指数对数函数不等式3} $e^x+e^y-2<e^{x+y}-1,\forall x,y>0.$

\item \label{proposition:指数对数函数不等式4} $e^{-x}\leqslant \frac{1}{1+x}\leqslant e^{x^2-x},\forall x>0.$
\end{enumerate}
\end{proposition}
\begin{proof}
\begin{enumerate}[(1)]
\item 令\(f(x)=\ln(1 + x)-\frac{x}{\sqrt{1 + x}}\)，\(x\geqslant 0\)，则
\begin{align*}
f'\left( x \right) =\frac{2\sqrt{1+x}-x-2}{2\left( 1+x \right) ^{\frac{3}{2}}}=-\frac{1+x-2\sqrt{1+x}+1}{2\left( 1+x \right) ^{\frac{3}{2}}}=-\frac{\left( \sqrt{1+x}-1 \right) ^2}{2\left( 1+x \right) ^{\frac{3}{2}}}<0,\forall x>0.
\end{align*}
故\(f\)在\((0,+\infty)\)上严格单调递减，又\(f\in C[0,+\infty)\)，因此\(f\)在\([0,+\infty)\)上也严格单调递减。从而
\[
f(x)\leqslant f(0)=0,\forall x>0.
\]
即\(\ln(1 + x)<\frac{x}{\sqrt{1 + x}},x>0\).

\item 注意到
\[
(e^x-1)(x^y-1)>0,\forall x,y>0,
\]
故
\begin{align*}
e^x+e^y<e^{x+y}+1\Longrightarrow e^x+e^y-2<e^{x+y}-1,\forall x,y>0.
\end{align*}

\item 由$e^x\geqslant 1+x,\forall x>0$可得
\begin{align*}
e^{-x}&=\frac{1}{e^x}\leqslant \frac{1}{1+y},\forall x>0.
\\
e^{x^2-x}&\geqslant 1+x^2-x=\frac{1+x^3}{1+x}\geqslant \frac{1}{1+x},\forall x>0.
\end{align*}
\end{enumerate}

\end{proof}

\begin{proposition}[经典初等不等式]\label{proposition:经典初等不等式}
\begin{enumerate}[(1)]
\item\label{proposition:常用不等式222-1} $\left| x^{\alpha}-y^{\alpha} \right|\leqslant \left| x-y \right|^{\alpha},\quad \forall x,y\geqslant 0,\,\,0<\alpha \leqslant 1.$

\item 设\(a,b\geqslant0\),则
\begin{align}\label{equation-12.777}
\begin{cases}
a^p + b^p\leqslant(a + b)^p\leqslant2^{p - 1}(a^p + b^p),& p\geqslant1,p\leqslant0\\
a^p + b^p\geqslant(a + b)^p\geqslant2^{p - 1}(a^p + b^p),& 0 < p < 1
\end{cases}
\end{align}
\end{enumerate}
\end{proposition}
\begin{note}
不等式左右是奇次对称的，我们可以设\(t = \frac{a}{b}\in[0,1]\)，于是\eqref{equation-12.777}两边同时除以$b^p$得
\[
\begin{cases}
t^p + 1\leqslant(t + 1)^p\leqslant2^{p - 1}(t^p + 1),& p\geqslant1,p\leqslant0\\
t^p + 1\geqslant(t + 1)^p\geqslant2^{p - 1}(t^p + 1),& 0 < p < 1
\end{cases}.
\]
\end{note}
\begin{proof}
\begin{enumerate}[(1)]
\item 不妨设$x\geqslant y>0$,令$t=\frac{x}{y}$,$u=t-1$,则
\[
\left| x^{\alpha}-y^{\alpha} \right|\leqslant \left| x-y \right|^{\alpha}\Longleftrightarrow t^{\alpha}-1\leqslant \left( t-1 \right)^{\alpha}\Longleftrightarrow \left( u+1 \right)^{\alpha}\leqslant u^{\alpha}+1.
\]
再由\hyperref[theorem:Bernoulli不等式]{Bernoulli不等式}知上式成立.

\item 考虑\(f(t)\triangleq\frac{(t + 1)^p}{1 + t^p},t\in[0,1]\)，我们有
\[
f'(t)=p(t + 1)^{p - 1}\frac{1 - t^{p - 1}}{(1 + t^p)^2}
\begin{cases}
\geqslant0,& p\geqslant1,p\leqslant0\\
<0,& 0 < p < 1
\end{cases}
\]
于是
\[
\begin{cases}
2^{p - 1}=f(1)\geqslant f(t)\geqslant f(0)=1,& p\geqslant1,p\leqslant0\\
2^{p - 1}=f(1)\leqslant f(t)\leqslant f(0)=1,& 0 < p < 1
\end{cases}
\]
这就完成了证明.
\end{enumerate}

\end{proof}
























































































































\end{document}